%%%%%%%%%%%%%%%%%%%%%%%%%%%%%%%%%%%%%%%%%%%%%%%%%%%%%%%%%%
%   Autoren:
%   Prof. Dr. Bernhard Drabant
%   Prof. Dr. Dennis Pfisterer
%   Prof. Dr. Julian Reichwald
%%%%%%%%%%%%%%%%%%%%%%%%%%%%%%%%%%%%%%%%%%%%%%%%%%%%%%%%%%

%%%%%%%%%%%%%%%%%%%%%%%%%%%%%%%%%%%%%%%%%%%%%%%%%%%%%%%%%%
%	ANLEITUNG: 
%   1. Ersetzen Sie firmenlogo.jpg im Verzeichnis img
%   2. Passen Sie alle Stellen im Dokument an, die mit 
%      @stud 
%      markiert sind
%%%%%%%%%%%%%%%%%%%%%%%%%%%%%%%%%%%%%%%%%%%%%%%%%%%%%%%%%%

%%%%%%%%%%%%%%%%%%%%%%%%%%%%%%%%%%%%%%%%%%%%%%%%%%%%%%%%%%
%	ACHTUNG: 
%   Für das Erstellen des Literaturverzeichnisses wird das 
%   modernere Paket biblatex in Kombination mit biber 
%   verwendet - nicht mehr das ältere Paket BibTex!
%
%   Bitte stellen Sie Ihre TeX-Umgebung entsprechend ein (z.B. TeXStudio): 
%   Einstellungen --> Erzeugen --> Standard Bibliographieprogramm: biber
%%%%%%%%%%%%%%%%%%%%%%%%%%%%%%%%%%%%%%%%%%%%%%%%%%%%%%%%%%

\documentclass[fontsize=12pt,BCOR=5mm,DIV=12,parskip=half,listof=totoc,
               paper=a4,toc=bibliography,pointlessnumbers]{scrreprt}
               
\usepackage[utf8]{inputenc}

%% LANGUAGE SETTINGS
%
% @stud: Sprache ggf. anpassen
%
%\usepackage[ngerman]{babel} 	        % german language
%\usepackage[german=quotes]{csquotes} 	% correct quoting using \enquote{}
\usepackage[english]{babel}          % english language
\usepackage{csquotes} 	              % correct quoting using \enquote{}

%%%%%%%%%%%%%
%% ZITIERSTIL
%%%%%%%%%%%%%
%
% @stud: Zitierstil in package biblatex unten wählen
%
% NUMERIC Style - e. g. [12]
% style=numeric 
%
% IEEE Style - numeric kind of style 
% style=ieee 
%
% ALPHABETIC Style - e. g. [AB12]
% style=alphabetic 
%
% HARVARD Style 
% style=apa 
%
% CHICAGO Style 
% style=authoryear
%
% Position des Zitats:
%
% autocite=inline 
%
% (!!) FOOTNOTE POSITION NOT RECOMMENDED IN MINT DOMAIN:
% autocite=footnote
%
\usepackage[backend=biber, autocite=inline, style=numeric]{biblatex} 	
\usepackage{makeidx}                  % allows index generation
\usepackage{listings}	                %Format Listings properly
\usepackage{lipsum}                   % Blindtext
\usepackage{graphicx}                 % use various graphics formats
\usepackage[german]{varioref}         % nicer references \vref
\usepackage{caption}	                % better Captions
\usepackage{booktabs}                 % nicer Tabs
\usepackage[hidelinks=true]{hyperref} % keine roten Markierungen bei Links
\usepackage{fnpct}                    % Correct superscripts 
\usepackage{calc}                     % Used for extra space below footsepline, in particular
\usepackage{array}
\usepackage{acronym}
\usepackage{algorithm}
\usepackage{algpseudocode}
\usepackage{setspace}
\usepackage{tocloft}
\usepackage[T1]{fontenc}
\usepackage{tabularx}
\usepackage[utf8]{inputenc}

% Definitionen und Commands
\newcommand{\indextype}{numeric}
\newcommand{\abs}{\par\vskip 0.2cm\goodbreak\noindent}
\newcommand{\nl}{\par\noindent}
\newcommand{\mcl}[1]{\mathcal{#1}}
\newcommand{\nowrite}[1]{}
\newcommand{\NN}{{\mathbb N}}
\newcommand{\imagedir}{img}
\newcommand{\TitelDerArbeit}[1]{\def\DerTitelDerArbeit{#1}\hypersetup{pdftitle={#1}}}
\newcommand{\Dozent}[1]{\def\DerDozent{#1}}
\newcommand{\AutorDerArbeit}[1]{\def\DerAutorDerArbeit{#1}\hypersetup{pdfauthor={#1}}}
\newcommand{\Firma}[1]{\def\DerNameDerFirma{#1}}
\newcommand{\Kurs}[1]{\def\DieKursbezeichnung{#1}}
\newcommand{\Abteilung}[1]{\def\DerNameDerAbteilung{#1}}
\newcommand{\Studiengangsleiter}[1]{\def\DerStudiengangsleiter{#1}}
\newcommand{\WissBetreuer}[1]{\def\DerWissBetreuer{#1}}
\newcommand{\FirmenBetreuer}[1]{\def\DerFirmenBetreuer{#1}}
\newcommand{\Bearbeitungszeitraum}[1]{\def\DerBearbeitungszeitraum{#1}}
\newcommand{\Abgabedatum}[1]{\def\DasAbgabedatum{#1}}
\newcommand{\MatrikelnummerEins}[1]{\def\DieMatrikelnummerEins{#1}}
\newcommand{\MatrikelnummerZwei}[1]{\def\DieMatrikelnummerZwei{#1}}
\newcommand{\MatrikelnummerDrei}[1]{\def\DieMatrikelnummerDrei{#1}}
\newcommand{\Studienrichtung}[1]{\def\DieStudienrichtung{#1}}
\newcommand{\ArtDerArbeit}[1]{\def\DieArtDerArbeit{#1}}
\newcommand{\Literaturverzeichnis}{Literaturverzeichnis}

% Page Layout
\oddsidemargin=0mm
\evensidemargin=0mm
\textwidth=159mm
\topmargin=-18mm
\headsep=10mm
\textheight=251mm
\footheight=15mm

\makeindex

%%%%%%%%%%%%%%%%%%%%%%%%%%%%%%%%%%%
% LITERATURVERZEICHNIS
% @stud: Literaturverzeichnis in Datei bibliography.bib anpassen. 
%
% [Alternative zu Verwendung von \initializeBibliography: Citavi ... (dazu eigenes LaTex Coding verwenden)]
%
\addbibresource{bibliography.bib}
\DefineBibliographyStrings{ngerman}{andothers = {{et\,al\adddot}},}

% Elementare Konfigurationen und Definitionen werden geladen 
% @stud: gegebenenfalls anpassen
%
% !TEX root =  master.tex

%%%%%%%%%%%%%%%%%%%%%%%%%%%%%%%%%%%%%%%%%%%%%%%%%%%%%%%%%%%%%%%%%%
%	ANLEITUNG: 
% Passen Sie gegebenenfalls alle Stellen im Dokument an, die mit 
% @stud 
% markiert sind.
%%%%%%%%%%%%%%%%%%%%%%%%%%%%%%%%%%%%%%%%%%%%%%%%%%%%%%%%%%%%%%%%%%

%%
%% @stud
%%
%% LANGUAGE SETTINGS
\usepackage[english]{babel}          % english language
\usepackage{csquotes} 	              % correct quoting using \enquote{}
\usepackage{makeidx}                  % allows index generation
\usepackage{listings}	                % Format Listings properly
\usepackage{lipsum}                   % Blindtext
\usepackage{graphicx}                 % use various graphics formats
\usepackage[german]{varioref}         % nicer references \vref
\usepackage{caption}	                % better Captions
\usepackage{booktabs}                 % nicer Tabs
\usepackage[hidelinks=false]{hyperref} % keine roten Markierungen bei Links
\usepackage{fnpct}                    % Correct superscripts 
\usepackage{calc}                     % Used for extra space below footsepline, in particular
\usepackage{array}
\usepackage{acronym}
\usepackage{algorithm}
\usepackage{algpseudocode}
\usepackage{setspace}
\usepackage{tocloft}

%% Schriftarten- und Zeichenpakete
\usepackage[T1]{fontenc}
\usepackage[utf8]{inputenc}

%%
%% @stud
%%
%%	FONT SELECTION: Schriftarten und Schriftfamilie
%%%%%%%%%%%%%
%% SCHRIFTART
%%%%%%%%%%%%%
% 0) without decomment: normal font families 
% ...
% 1) Latin Modern 
%\usepackage{lmodern}        
% 2) Times 
%\usepackage{mathptmx}         
% 3) Helvetica
%\usepackage[scaled=.92]{helvet} 
%%%%%%%%%%%%%%%%%%
%%	SCHRIFTFAMILIE
%%%%%%%%%%%%%%%%%%
% ohne Serifen
\renewcommand*{\familydefault}{\sfdefault}
\addtokomafont{disposition}{\sffamily}
%
% mit Serifen
%\renewcommand*{\familydefault}{\rmdefault}
%\addtokomafont{disposition}{\rmfamily}
%
% Typewriter
%\renewcommand*{\familydefault}{\ttdefault}
%\addtokomafont{disposition}{\ttfamily}

%%
%% @stud
%%
%% Uncomment the following lines to support hard URL breaks in bibliography 
%\apptocmd{\UrlBreaks}{\do\f\do\m}{}{}
%\setcounter{biburllcpenalty}{9000}% Kleinbuchstaben
%\setcounter{biburlucpenalty}{9000}% Großbuchstaben

%%
%% @stud
%%
%% FOOTNOTES: Count footnotes over chapters
%% \counterwithout{footnote}{chapter}

%	ACRONYMS
\makeatletter
\@ifpackagelater{acronym}{2015/03/20}
{\renewcommand*{\aclabelfont}[1]{\textbf{{\acsfont{#1}}}}}{}
\makeatother

%	LISTINGS
% @stud: ggf. Namen/Text anpassen (englisch)
\renewcommand{\lstlistingname}{Quelltext} 
\renewcommand{\lstlistlistingname}{Quelltextverzeichnis}
\lstset{numbers=left,
	numberstyle=\tiny,
	captionpos=b,
	basicstyle=\ttfamily\small}

%	ALGORITHMS
% @stud: ggf. Namen/Text anpassen (englisch)
\renewcommand{\listalgorithmname}{Algorithmenverzeichnis}
\floatname{algorithm}{Algorithmus}

%	PAGE HEADER / FOOTER
%	Warning: There are some redefinitions throughout the master.tex-file!  DON'T CHANGE THESE REDEFINITIONS!
\RequirePackage[automark]{scrlayer-scrpage}
%alternatively with separation lines: \RequirePackage[automark,headsepline,footsepline]{scrlayer-scrpage}

\renewcommand{\chaptermarkformat}{}
\RedeclareSectionCommand[beforeskip=0pt]{chapter}
\clearpairofpagestyles

%\ifoot[\rule{0pt}{\ht\strutbox+\dp\strutbox}DHBW Mannheim]{\rule{0pt}{\ht\strutbox+\dp\strutbox}DHBW Mannheim}
\ofoot[\rule{0pt}{\ht\strutbox+\dp\strutbox}\pagemark]{\rule{0pt}{\ht\strutbox+\dp\strutbox}\pagemark}
\ohead{\headmark}

%\newcommand{\TitelDerArbeit}[1]{\def\DerTitelDerArbeit{#1}\hypersetup{pdftitle={#1}}}
%\newcommand{\AutorDerArbeit}[1]{\def\DerAutorDerArbeit{#1}\hypersetup{pdfauthor={#1}}}
%\newcommand{\Firma}[1]{\def\DerNameDerFirma{#1}}
%\newcommand{\Kurs}[1]{\def\DieKursbezeichnung{#1}}
%\newcommand{\Abteilung}[1]{\def\DerNameDerAbteilung{#1}}
%\newcommand{\Studiengangsleiter}[1]{\def\DerStudiengangsleiter{#1}}
%\newcommand{\WissBetreuer}[1]{\def\DerWissBetreuer{#1}}
%\newcommand{\FirmenBetreuer}[1]{\def\DerFirmenBetreuer{#1}}
%\newcommand{\Bearbeitungszeitraum}[1]{\def\DerBearbeitungszeitraum{#1}}
%\newcommand{\Abgabedatum}[1]{\def\DasAbgabedatum{#1}}
%\newcommand{\Matrikelnummer}[1]{\def\DieMatrikelnummer{#1}}
%\newcommand{\Studienrichtung}[1]{\def\DieStudienrichtung{#1}}
%\newcommand{\ArtDerArbeit}[1]{\def\DieArtDerArbeit{#1}}
%\newcommand{\Literaturverzeichnis}{Literaturverzeichnis}

\newcommand{\settingBibFootnoteCite}{
	\setlength{\bibparsep}{\parskip}		  % Add some space between biblatex entries in the bibliography
	\addbibresource{bibliography.bib}	    % Add file bibliography.bib as biblatex resource
	\DefineBibliographyStrings{ngerman}{andothers = {{et\,al\adddot}},}
}

\newcommand{\setTitlepage}{
	% !TEX root =  master.tex
% @stud: ggf. Namen/Text anpassen (englisch)
\begin{titlepage}
\begin{minipage}{\textwidth}
		\vspace{-2cm}
		\noindent \hfill \includegraphics{\imagedir/logo.jpg}
\end{minipage}
\vspace{1em}
%\sffamily
\begin{center}
	{\textsf{\large Duale Hochschule Baden-W\"urttemberg Mannheim}}\\[4em]
	{\textsf{\textbf{\large{\DieArtDerArbeit}arbeit}}}\\[6mm]
	{\textsf{\textbf{\Large{}\DerTitelDerArbeit}}} \\[1.5cm]
	{\textsf{\textbf{\large{}Studiengang Informatik}}\\[6mm]
	\textsf{\textbf{Studienrichtung \DieStudienrichtung}}}\vspace{10em}
	
	\begin{minipage}{\textwidth}
		\begin{tabbing} 
		Wissenschaftliche(r) Betreuer(in): \hspace{0.85cm}\=\kill
		Verfasser(in): \> \DerAutorDerArbeit \\[1.5mm] 
		Matrikelnummer: \> \DieMatrikelnummerEins \\[1.5mm]
		Matrikelnummer: \> \DieMatrikelnummerZwei \\[1.5mm]
		Matrikelnummer: \> \DieMatrikelnummerDrei \\[1.5mm]
		%Firma: \> \DerNameDerFirma  \\[1.5mm]
		%Abteilung: \> \DerNameDerAbteilung \\[1.5mm]
		Kurs: \> \DieKursbezeichnung \\[1.5mm]
		Dozent: \> \DerDozent\\[1.5mm]
		%Studiengangsleiter: \> \DerStudiengangsleiter \\[1.5mm]
		%Wissenschaftliche(r) Betreuer(in): \> \DerWissBetreuer \\[1.5mm]
		%Firmenbetreuer(in): \> \DerFirmenBetreuer \\[1.5mm]
		Bearbeitungszeitraum: \> \DerBearbeitungszeitraum\\[1.5mm]
		
%		alternativ:\\[1.5mm]
%		Eingereicht: \> \DasAbgabedatum	
		\end{tabbing}
	\end{minipage}
\end{center}
\end{titlepage}
	\pagenumbering{roman} % Römische Seitennummerierung
	\normalfont	
}

\newcommand{\initializeText}{
	\clearpage
	\ihead{\chaptername~\thechapter} % Neue Header-Definition
	\pagenumbering{arabic}           % Arabische Seitenzahlen
}

\newcommand{\initializeBibliography}{
	\ihead{}
	\printbibliography[title=\Literaturverzeichnis] 
	\cleardoublepage
}

\newcommand{\initializeAppendix}{
	\appendix
  \ihead{}
  \cftaddtitleline{toc}{chapter}{Anhang}{}
}



% @stud
%
% PERSÖNLICHE ANGABEN (BITTE VOLLSTÄNDIG EINGEBEN zwischen den Klammern: {...})
%
\ArtDerArbeit{Projektbericht} % "Bachelor" oder "Projekt" wählen
\TitelDerArbeit{Drei-Schicht Anwendung}
\AutorDerArbeit{Moritz Werr, Phil Richter, Max Stege}
\Abteilung{<Ihre Abteilung>}
\Firma{<Ihre Firma>}
\Kurs{TINF22AI2}
\Studienrichtung{Angewandte Informatik}
\MatrikelnummerEins{5401527}
\MatrikelnummerZwei{4164342}
\MatrikelnummerDrei{7285772}
\Studiengangsleiter{<Ihr Studiengangsleiter>}
\WissBetreuer{<Ihr(e) wissenschaftliche(r) Betreuer(in)>}
\FirmenBetreuer{<Ihr(e) Firmenbetreuer(in)>}
\Dozent{Jochen Kluger}
\Bearbeitungszeitraum{01.10.2024 -- 20.12.2024}
\Abgabedatum{dd.mm.yyyy}

\begin{document}

\setTitlepage

%%%%%%%%%%%%%%%%%%%%%%%%%%%%%%%%%%%
% EHRENWÖRTLICHE ERKLÄRUNG
%
% @stud: ewerkl.tex bearbeiten
%
% !TEX root =  master.tex
\clearpage
\chapter*{Ehrenwörtliche Erklärung}

% Wird die folgende Zeile auskommentiert, erscheint die ehrenwörtliche
% Erklärung im Inhaltsverzeichnis.

% \addcontentsline{toc}{chapter}{Ehrenwörtliche Erklärung}
Ich versichere hiermit, dass ich die vorliegende Arbeit mit dem Titel ``\textit{\DerTitelDerArbeit}'' selbstständig verfasst und 
keine anderen als die angegebenen Quellen und Hilfsmittel benutzt habe. Ich versichere zudem, dass die eingereichte elektronische 
Fassung mit der gedruckten Fassung übereinstimmt.

\vspace{3cm}
%Ort, Datum \hfill \DerAutorDerArbeit
 
\cleardoublepage  
%%%%%%%%%%%%%%%%%%%%%%%%%%%%%%%%%%%

%%%%%%%%%%%%%%%%%%%%%%%%%%%%%%%%%%%
% SPERRVERMERK
%
% @stud: nondisclosurenotice.tex bearbeiten
%
%%%%%%%%%%%%%%%%%%%%%%%%%%%%%%%%%%%

%%%%%%%%%%%%%%%%%%%%%%%%%%%%%%%%%%%
%	KURZFASSUNG
%
% @stud: acknowledge.tex bearbeiten
%
%%%%%%%%%%%%%%%%%%%%%%%%%%%%%%%%%%%

%%%%%%%%%%%%%%%%%%%%%%%%%%%%%%%%%%%
% VERZEICHNISSE und ABSTRACT
%
% @stud: ggf. nicht benötigte Verzeichnisse auskommentieren/löschen
%
%	Listingsverzeichnis / Quelltextverzeichnis
\lstlistoflistings
\cleardoublepage
\onehalfspacing
\initializeText

%%%%%%%%%%%%%%%%%%%%%%%%%%%%%%%%%%%%%%%%%%%%%%%%%%%%%%%%%%%%%%%%%%%%%%%%%%%%%%%%%%%%%%%%%%
% KAPITEL UND ANHÄNGE
%
% @stud:
%   - nicht benötigte: auskommentieren/löschen
%   - neue: bei Bedarf hinzufügen mittels input-Kommando an entsprechender Stelle einfügen
%%%%%%%%%%%%%%%%%%%%%%%%%%%%%%%%%%%%%%%%%%%%%%%%%%%%%%%%%%%%%%%%%%%%%%%%%%%%%%%%%%%%%%%%%%

%%%%%%%%%%%%%%%%%%%%%%%%%%%%%%%%%%%
% KAPITEL
%
% @stud: einzelne Kapitel bearbeiten und eigene Kapitel hier einfügen
%
% Einleitung
%% !TEX root =  master.tex
\chapter{Einleitung (Moritz Werr)}

\nocite{*}

Im Rahmen dieses Projektes wurde eine Anwendung in drei Schichten aufgeteilt mithilfe von Docker und Docker-Compose. Diese wurde im Rahmen eines vorherigen Semesters implementiert als eine monolithische Serveranwendung, obwohl diese jedoch eigentlich aus mehreren Komponenten besteht. Generell besteht diese aus drei Teilen: dem Webserver und der Programmierlogik für die Anwendung, einem Datenbankserver für die persistente Speicherung der Daten und einer TLS-Verschlüsselung für die Sicherheit der Datenübertragung. Vorher liefen diese alle auf einem Server, es war eine monolithische Architektur. Nun wird diese in ihren einzelnen Funktionalitäten aufgespalten, hin zu einer Mikroservicearchitektur. Jede einzelne Funktionalität wird in einen eigenen Softwarecontainer übertragen und laufen isoliert auf der gleichen Maschine, indem sie sich den Betriebsystemkernel teilen \cite[Vgl. S. 152444]{watada_emerging_2019}. Dieser Prozess nennt sich Containerisierung.

Durch die Umwandlung in eine Mikroservicearchitektur wird ein hoher Grad an Flexibilität gewonnen. Die einzelnen Komponenten arbeiten unabhängig voneinander und ergeben gemeinsam ein funktionierendes System \cite[Vgl. S.1f]{fowler_microservices_2015}. Dafür muss die Anwendung jedoch in manchen Fällen  umgeschrieben werden. Durch den Umbau zu einer Mikroservicearchitektur, wird oftmals die Komplexität der Anwendung erhöht \cite[Vgl. S.10]{su_modular_2024}, weswegen noch weitere Werkzeuge zur Verwaltung der Anwendung benötigt werden. Hilfswerkzeuge zur Verwaltung und Konfiguration der Anwendung werden benötigt. 

Weiterhin sind nicht nur die einzelnen Komponenten austauschbar, sondern können nun auch weitere Funktionalitäten durch die Containerisierung erreicht werden. Es können sehr einfach Healthchecks implementiert werden, die üverprüfen, ob einzelne Komponenten noch richtig funktionieren \cite{noauthor_services_0100}. Wenn dies nicht mehr der Fall ist, können die einzelnen Komponenten der drei Schichten neugestartet werden \cite{noauthor_services_0100}. Dafür muss definiert werden wie diese miteinander zusammenarbeiten. Wenn nun ein einzelner Container Probleme hat, wird dies erkannt und er kann automatisch neugestartet werden. Weiterhin sorgt dies auch dafür, dass Container in der richtigen Reihenfolge gestartet werden \cite{noauthor_services_0100}.

Im Rahmen dieses Projektes wurden Docker und Docker-Compose verwendet aufgrund ihrer guten Dokumentation und großen Anzahl an Features. Kernbestandteil war Docker zur Containerisierung und Docker-Compose zur Verwaltung der Container. Es musste ein neues Dockerfile für den Webserver erstellt werden, zwei vorhandene Containerimages wurden für die TLS-Verschlüsselung (traefik) bzw. die Datenspeicherung (mariadb) verwendet und ein Initialisierungsskript für die Datenbank wurde erstellt.

% mehrere Grundlagen- und Forschungs-Kapitel
%% !TEX root =  master.tex
\chapter{Docker als Softwarecontainer (Moritz Werr)}

Für die Containerisierung wurde Docker verwendet. Softwarecontainer bieten viele Vorteile gegenüber traditioneller Virtualisierung. Traditionelle Virtualisierung wird verwendet um Ressourcen effizienter zu nutzen. Dabei werden viele schwache Maschinen in einer großen Maschine konsolidiert \cite[Vgl. S.152443]{watada_emerging_2019}. Diese teilen sich alle die gleichen Hardwareressourcen, welche über den Hypervisor auf die einzelnen virtuellen Maschinen verteilt wird. Allerdings teilen sich virtuelle Maschinen nicht ihre Softwarekomponenten. Dadurch wird oftmals der gleiche Programmcode öfters im Arbeitsspeicher vorkommen. Softwarecontainer funktionieren in dieser Hinsicht anders. Sie teilen sich normalerweise den Kernel mit dem Containerhost \cite[Vgl. S.152444]{watada_emerging_2019}. Diese sind dadurch um einiges leichtgewichtiger und schneller als traditionelle Virtualisierung \cite[Vgl. S.137]{fava_assessing_2024}. 
Aufgrund dieser Leichtgewichtigkeit und Geschwindigkeit lassen sich Docker-Anwendungen mit Orchestrierungswerkzeugen wie Kubernetes auch sehr leicht horizontal und vertikal skalieren \cite[Vgl. S.152458]{watada_emerging_2019}. 

Weiterhin sind Anwendungen innerhalb von Docker-Containern auch sehr portabel, da ein Container nicht vom Hostsystem abhängig ist, sondern von dessen Docker-Image. Alle Abhängigkeiten (Dependencies) sind in einem einzelnen Image gebündelt  \cite[Vgl. S.137]{fava_assessing_2024}. Dadurch sind containerisierte Anwendungen hoch portabel. Sie können auf jeder Maschine laufen, die eine kompatible Docker-Engine hat. Nicht nur läuft die Anwendung auf sehr vielen Maschinen, sondern es läuft auch überall die gleiche Version. Das Image sorgt für einen hohen Grad an Konsistenz und Wiederholbarkeit.

Außerhalb der hohen Effizienz und Portabilität von Softwarecontainern, bieten diese auch noch einen gewissen Grad an Sicherheit. Obwohl virtuelle Maschinen einen höheren Grad an Sicherheit bieten durch getrennte Kernel \cite[Vgl. S.152455]{watada_emerging_2019}, bieten Container einen guten Schutz durch Prozessisolierung, Dateisystemisolierung, Ressourcenlimitierung und Netzwerkisolation \cite[Vgl. S.152444ff]{watada_emerging_2019}. Dafür werden Funktionalitäten des Linux Kernels verwendet. Primär wären dies CGroups, Namespaces und MAC-Systemen (Mandatory Access Control) wie SELinux oder Apparmor. Cgroup sorgen dabei für die Resourcenlimitierung, Namespaces für die Prozess- und Netzwerkisolation und Apparmor für die Dateisystemisolierung.

Durch die Nutzung von Docker ist die Anwendung sehr performant, portabel und auch noch gut abgesichert. 
%% !TEX root =  master.tex
\chapter{Säubern der Daten (Moritz Werr)} \label{chap:säubern}


Nun ist die CSV vorbereitet für die eigentliche Säuberung. Die Daten liegen in den entsprechenden Spalten vor und alle uninteressanten Daten wurden entfernt.
Bevor die Daten jedoch in einer Datenbank in Tabellenform gespeichert werden können, sollten noch entsprechende Spalten bereinigt werden, dass nur noch Zahlenwerte ohne Einheiten dort aufzufinden sind.
Weiterhin beinhalten viele Spaltennamen immer noch Sonderzeichen und Leerzeichen, welche es schwierig macht die Spalten zu lesen.

Um diese Probleme zu addressieren, wurde R zur Bereinigung der Daten ausgewählt.
R ist eine Programmiersprache für statistische Berechnungen \cite{noauthor_r_nodate}. Es bietet mächtige Werkzeuge zur Auswertung und Speicherung von Daten \cite{noauthor_r_nodate}. Damit ist es eine exzellente Wahl um dieses Problem zu lösen.
Es bietet Möglichkeiten die Spaltennamen automatisch mit der Funktion  \lstinline |clean_names()| aus der Bibliothek \lstinline|janitor| umzubenennen \cite{noauthor_janitor_nodate}, indem es problematische Zeichen, wie z.B. Leerzeichen und Sonderzeichen, erkennt und durch unproblematische Zeichen ersetzt.
Weiterhin werden auch alle Großbuchstaben durch Kleinbuchstaben ersetzt.
Darüber hinaus ist R eine dynamisch typisierte Programmiersprache \cite{ewendorf_algorithmic_2019}. Datentypen werden je nach Kontext automatisiert erkannt.
Wenn die entsprechenden Spalten von allen nicht-numerischen Zeichen gesäubert werden, lassen sich diese problemlos in Zahlen umwandeln.
Die entsprechenden Einheiten werden nach der Spaltensäuberung stattdessen noch in den Spaltennamen hinzugefügt.

Wenn alle Spalten und Spaltennamen gesäubert bzw. angepasst wurden, können die Daten in einer Datenbank gespeichert werden.

Mit den vorher genannten Überlegungen im Hinterkopf ist folgendes Skript entstanden: 


\lstset{
	breaklines=true,         % Enable line wrapping
	breakatwhitespace=false, % Allow breaks at any character (not just whitespace)
	basicstyle=\ttfamily,    % Use monospaced font
}

\begin{lstlisting}[caption={\texttt{database.R}},captionpos=b]
#!/usr/bin/Rscript

library(RPostgres)

intel <- read.csv("clean.csv")		#Reading
intel <- janitor::clean_names(intel)	#Cleaning the column names

#cleaning and mutating columns to numbers
intel$max_turbo_taktfrequenz <- as.numeric(gsub(" GHz", "", intel$max_turbo_taktfrequenz))
intel$lithographie <- gsub("Intel 7", "10 nm", intel$lithographie)
intel$lithographie <- as.numeric(gsub(" nm", "", intel$lithographie))
intel$intel_turbo_boost_technik_2_0_taktfrequenz <- as.numeric(gsub(" GHz", "", intel$intel_turbo_boost_technik_2_0_taktfrequenz))
intel$grundtaktfrequenz_des_prozessors <- gsub(" \\| ", ".", intel$grundtaktfrequenz_des_prozessors)
intel$grundtaktfrequenz_des_prozessors <- as.numeric(gsub(" GHz", "", intel$grundtaktfrequenz_des_prozessors))
intel$cache <- gsub(" MB Intel Smart Cache", "", intel$cache)
intel$cache <- as.numeric(gsub(" MB", "", intel$cache))
intel$bus_taktfrequenz <- as.numeric(gsub(" GT/s", "", intel$bus_taktfrequenz))
intel$verlustleistung_tdp <- as.numeric(gsub(" W", "", intel$verlustleistung_tdp))
intel$intel_turbo_boost_max_technology_3_0_frequency <- gsub(" \\| ", ".", intel$intel_turbo_boost_max_technology_3_0_frequency)
intel$intel_turbo_boost_max_technology_3_0_frequency <- as.numeric(gsub(" GHz", "", intel$intel_turbo_boost_max_technology_3_0_frequency))
intel$single_p_core_turbo_frequency <- as.numeric(gsub(" GHz", "", intel$single_p_core_turbo_frequency))
intel$single_e_core_turbo_frequency <- as.numeric(gsub(" GHz", "", intel$single_e_core_turbo_frequency))
intel$e_core_base_frequency <- gsub(" GHz", "", intel$e_core_base_frequency)
intel$e_core_base_frequency <- as.numeric(gsub("900 MHz", "0.9", intel$e_core_base_frequency))
intel$total_l2_cache <- as.numeric(gsub(" MB", "", intel$total_l2_cache))
intel$processor_base_power <- as.numeric(gsub(" W", "", intel$processor_base_power))
intel$maximum_turbo_power <- as.numeric(gsub(" W", "", intel$maximum_turbo_power))
intel$grundtaktfrequenz_der_grafik <- as.numeric(gsub(" MHz", "", intel$grundtaktfrequenz_der_grafik))
intel$max_dynamische_grafikfrequenz <- as.numeric(gsub(" GHz", "", intel$max_dynamische_grafikfrequenz))
intel$max_videospeicher_der_grafik <- as.numeric(gsub(" GB", "", intel$max_videospeicher_der_grafik))
intel$x4k_unterstutzung <- gsub("Hz", "", intel$x4k_unterstutzung)
intel$x4k_unterstutzung <- as.numeric(gsub("Yes \\|  at ", "", intel$x4k_unterstutzung))


#Putting the units back into the column names
colnames(intel)[colnames(intel) == "max_turbo_taktfrequenz"] <- "max_turbo_taktfrequenz_GHz"
colnames(intel)[colnames(intel) == "lithographie"] <- "litographie_nm"
colnames(intel)[colnames(intel) == "intel_turbo_boost_technik_2_0_taktfrequenz"] <- "intel_turbo_boost_technik_2_0_taktfrequenz_GHz"
colnames(intel)[colnames(intel) == "grundtaktfrequenz_des_prozessors"] <- "grundtaktfrequenz_des_prozessors_GHz"
colnames(intel)[colnames(intel) == "cache"] <- "cache_MB"
colnames(intel)[colnames(intel) == "bus_taktfrequenz"] <- "bus_taktfrequenz_GT_per_s"
colnames(intel)[colnames(intel) == "verlustleistung_tdp"] <- "verlustleistung_tdp_W"
colnames(intel)[colnames(intel) == "intel_turbo_boost_max_technology_3_0_frequency"] <- "intel_turbo_boost_max_technology_3_0_frequency_GHz"
colnames(intel)[colnames(intel) == "single_p_core_turbo_frequency"] <- "single_p_core_turbo_frequency_GHz"
colnames(intel)[colnames(intel) == "single_e_core_turbo_frequency"] <- "single_e_core_turbo_frequency_GHz"
colnames(intel)[colnames(intel) == "e_core_base_frequency"] <- "e_core_base_frequency_GHz"
colnames(intel)[colnames(intel) == "total_l2_cache"] <- "total_l2_cache_MB"
colnames(intel)[colnames(intel) == "processor_base_power"] <- "processor_base_power_W"
colnames(intel)[colnames(intel) == "maximum_turbo_power"] <- "maximum_turbo_power_W"
colnames(intel)[colnames(intel) == "grundtaktfrequenz_der_grafik"] <- "grundtaktfrequenz_der_grafik_MHz"
colnames(intel)[colnames(intel) == "max_dynamische_grafikfrequenz"] <- "max_dynamische_grafikfrequenz_GHz"
colnames(intel)[colnames(intel) == "max_videospeicher_der_grafik"] <- "max_videospeicher_der_grafik_GB"
colnames(intel)[colnames(intel) == "x4k_unterstutzung"] <- "x4k_unterstutzung_at"


con <- dbConnect(
RPostgres::Postgres(),
dbname = "bda",
host = "localhost" ,
port = 5432,
user = "bda",
password = "bda",
)

dbWriteTable(con, "intel", intel, row.names = FALSE, overwrite = TRUE)
dbDisconnect(con)
\end{lstlisting}

Als erstes wird im Skript die transponierte CSV eingelesen und die Spaltennamen mit Hilfe der  \lstinline |janitor| Library gesäubert. Danach kommt ein großer Block an Funktionen, die die einzelnen Spalten von Einheiten und komischer Formatierung säubern.
Danach kommt ein großer Block zur Umbenennung der Spaltennamen. 
Im letzten Schritt wird der Dataframe in einer existierenden Datenbank gespeichert. 
Ein großer Vorteil an der Funktion ist, dass die Tabelle innerhalb der Datenbank vorher nicht existieren muss. 
R erstellt von selber eine neue Tabelle mit entsprechenden Datentypen basierend auf den Datentypen innerhalb des Dataframes. Daher ist es wichtig die Datentypen innerhalb des Dataframes schon zu Nummern zu konvertieren.

%% !TEX root =  master.tex
\chapter{Speichern der Daten (Moritz Werr)}

Aus dem vorherigen Kapitel ist bereits hervorgegangen, dass die Daten in einer Datenbank gespeichert werden.
In diesem Fall ist die Wahl auf PostgreSQL gefallen.
PostgreSQL ist ein Open-Source, relationales Datenbankmanagementsystem \cite{noauthor_postgresql_nodate}. Es ist komplett ACID-kompatibel \cite{noauthor_postgresql_nodate} und speichert Daten in geordneter Form als Tabellen ab \cite{noauthor_51_2024}.
Da die Datenquelle bereits als CSV geliefert wird, bietet es sich an diese in einer entsprechenden Datenbank zu speichern.
CSV-Dateien speichern bereits Daten in Tabellenform, weswegen die Wahl direkt auf eine relationale SQL-Datenbank gefallen ist.

Es gibt viele Alternativen zu PostgreSQL, diese sind aber teils kostenpflichtig oder bieten die gleichen bzw. weniger Features wie PostgreSQL.
Proprietäre Datenbankmanagementsysteme wie Oracle RDBMS oder Microsoft SQL-Server sind für große, Enterprise-Skala Deployments konzipiert und entsprechend kommerzielle Lizenzen \cite{noauthor_microsoft_nodate}\cite{noauthor_oracle_nodate}.
Kostenlose Optionen sind daher bevorzugt.
MariaDB bzw. MySQL und SQLite wären Open-Source Alternativen zu PostgreSQL, welche auch gut geeignet sind\cite{noauthor_sqlite_nodate}\cite{noauthor_mysql_nodate}.
SQLite jedoch bietet keine Serverfunktionalität weswegen nur lokal gearbeitet werden kann\cite{noauthor_sqlite_nodate}.
MariaDB/MySQL haben keinen nennenswerten Funktionalitätsunterschied, aber wurden nicht gewählt, weil mit PostgreSQL mehr Erfahrung im Team existiert\cite{noauthor_mysql_nodate}\cite{noauthor_postgresql_nodate-1}.

Nachdem das Skript aus dem vorherigen Kapitel durchgelaufen ist, ist das Datenbankschema in Abbildung \ref{fig:datenbankschema} zu sehen.

\begin{figure}
	\centering 
	\includegraphics[width=\textwidth]{\imagedir/datenbank_schema.png} 
	\captionsetup{format=hang}
	\caption[Datenbankschema]{\label{fig:datenbankschema}Datenbankschema erstellt durch das R-Skript}
\end{figure}
%% !TEX root =  master.tex
\chapter{Die Webanwendung (Phil Richter)}
Die Webanwendung besteht aus einem Internetforum, in dem Benutzer Beiträge erstellen, kommentieren und liken
können. Um am Forum teilzunehmen, muss sich der Benutzer mit einer Emailadresse und einem Passwort registrieren und einloggen.
Falls der Benutzer sein Passwort vergisst, kann er dieses mit Hilfer einer Emailadresse zurücksetzen und ein neues vergeben. 
Des Weiteren existiert ein Admin Account mit erweiterten Rechten. Dieser kann Beiträge von allen Nutzern
löschen, während normale Benutzer nur ihre eigenen Beiträge löschen können. Zudem kann dieser sowohl die Emailadresse als auch
das Passwort eines Benutzers ändern. Falls notwendig kann ein Admin auch einen Account löschen.

Die Anwendung besteht im Frontend aus HTML \cite{noauthor_html_nodate}, CSS \cite{noauthor_cascading_nodate} und JavaScript
\cite{noauthor_learn_nodate}. Das Backend wurde mit PHP \cite{noauthor_php_2024} geschrieben. Zusätzlich existiert noch eine
Datenbank, in der die Benutzerdaten und Beiträge gespeichert werden. Die Datenbank wurde mit MariaDB \cite{noauthor_mariadb_nodate}
realisiert. Die Mails werden mit msmtp \cite{noauthor_msmtp_nodate} über einen SMTP-Server von Gmail versendet.

\chapter{Webserver}
Für den Webserver wurde ein Dockerfile implementiert welches wiefolgt aussieht:
\begin{lstlisting}[caption={\texttt{dockerfile}},captionpos=b, breaklines=true]
    FROM php:8.1-apache
    COPY src/ /var/www/html
    RUN apt-get update && apt-get install -y msmtp && rm -rf /var/cache/apt/lists
    COPY msmtprc /etc/msmtprc
    RUN chmod 600 /etc/msmtprc && chown www-data: /etc/msmtprc
    RUN mkdir /var/www/html/pictures
    RUN sed -i 's/AllowOverride None/AllowOverride All/g' /etc/apache2/apache2.conf
    RUN printf "AddType application/x-httpd-php .html\nAddType application/x-httpd-php .htm\n" >> /etc/apache2/apache2.conf
    RUN a2enmod rewrite
    RUN chown www-data: -R /var/www/html/
    RUN docker-php-ext-install mysqli && docker-php-ext-enable mysqli
    RUN echo "sendmail_path = /usr/bin/msmtp -t" >> /usr/local/etc/php/php.ini-development
    RUN cp /usr/local/etc/php/php.ini-development /usr/local/etc/php/php.ini
    CMD apachectl -DFOREGROUND
\end{lstlisting}

Im ersten Schritt wird in Zeile 1 das Basisimage \lstinline|php: 8.1-apache| gesetzt. Dies ist ein bereis vorkonfiguriertes
Image mit PHP 8.1 und Apache Webserver. In Zeile 2 wird der Quellcode der Webanwendung, welche im Ordner \lstinline|src/| liegt,
in den Ordner \lstinline|/var/www/html| des Containers kopiert. Damit die Webanwendung erfolgreich ausgefürt wird, müssen
verschiedene Konfigurationen von Apache vorgenommen werden. Dies passiert in Zeile 7 bis 9. Dabei wird unteranderem das Apache
Modul \lstinline|rewrite| aktiviert, sowie angepasst, dass die Dateiendungen \lstinline|.html| und \lstinline|.htm| als PHP-Dateien
interpretiert werden.


Zusätzlich werden weitere, für die Webanwendung notwendige Komponenten installiert. Dazu gehört für den Mailversand
\lstinline|msmtp| in Zeile 3 und für die Datenbank \lstinline|mysqli| in Zeile 11. Damit msmtp verwendet wird,
wird in Zeile 12 der \lstinline|sendmail_path| in der PHP-Konfiguration korrekt gesetzt. Zudem wird in Zeile 4
die Konfigurationsdatei \lstinline|msmtprc| in den Container kopiert. Diese Datei enthält die Konfigurationen wie unter
anderem Host und Port für den Mailversand.

Des Weitren wird in Zeile 6 ein Ordner \lstinline|pictures| erstellt, in welchem später die von den Nutzern hochgelandenen
Bilder gespeichert werden. In Zeile 5 und 10 wird der Besitzer der Dateien  und Ordner auf \lstinline|www-data| gesetzt.
Dadurch wird sichergestellt, dass der Apache Webserver auf die Dateien zugreifen kann.
Zum Schluss wird in Zeile 14 der Apache Webserver gestartet.


Um den Webserver als eigenständigen Container zu betreiben wird folgender Code in der \lstinline|docker-compose.yml| Datei
definiert:
\begin{lstlisting}[caption={\texttt{docker-compose.yml}},captionpos=b]
    webserver:
    image: webserver
    build: .
    volumes:
      - pictures:/var/www/html/pictures
    labels:
      - "traefik.enable=true"
      - "traefik.http.routers.webserver.rule=Host(`localhost`)"
      - "traefik.http.services.webserver.loadbalancer.server.port=80"
      - "traefik.http.routers.webserver.entrypoints=web,websecure"
      - "traefik.http.routers.webserver.tls=true"
    environment:
      DATABASE_USER: web_eng_user
      DATABASE_HOST: mariadb
      DATABASE_PASS: linux
      DATABASE_NAME: web_eng
      GMAIL_APP_PASSWORD: ${GMAIL_APP_PASSWORD}
    depends_on:
      db:
        condition: service_healthy
        restart: true
    healthcheck:
      test: ["CMD", "curl", "-f", "http://localhost"]
      interval: 10s
      timeout: 5s
      retries: 3
      start_period: 10s
\end{lstlisting}

Dabei wird in Zeile 2 das Image \lstinline|webserver| definiert, welches in Zeile 3 aus dem aktuellen Verzeichnis, dem dockerfile,
gebaut wird. In Zeile 5 wird ein Volume \lstinline|pictures| definiert. Dies ist dazu notwendig, dass die hochgeladenen Bilder
auch nach einem Neustart des Containers noch vorhanden sind. Die in Zeile 6 und folgenden definierten Labels sind dazu notwendig,
dass der Webserver über den TLS-Proxy \lstinline|traefik| erreichbar ist. Außerdem wird festgelegt, dass der Webserver auf
Localhost erreichbar ist und auf Port 80 lauscht. Zusätzlich wird tls aktiviert.

Um auf die Datenbank zugreifen zu können und Mails über Gmail versenden zu können, werden in Zeile 12 und folgende 
die notwendigen Umgebungsvariablen definiert und gesetzt. Dabei wird das Passwort für den Gmail Account aus einer
Lokalen .env Datei gelesen. Somit wird sichergestellt, dass das Passwort nicht im Klartext vorliegt.

Da der Webserver auf die Datenbank angewiesen ist, wird in Zeile 18 bis 21 eine Abhängigkeit definiert. Dabei
wird überprüft, ob die Datenbank erreichbar ist. Falls nicht, wird diese neu gestartet. Des Weiteren wird in Zeile 22 bis 27
ein Healthcheck definiert. Dabei wird alle 10 Sekunden mit Hilfe eines Curl-Befehls überprüft, ob der Webserver erreichbar ist.
Falls der Healthcheck dreimal fehlschlägt, gilt der Container als "unhealthy".

% Fazit und Ausblick
%% !TEX root =  master.tex
\chapter{Starten der Anwendung (Moritz Werr)}


%%%%%%%%%%%%%%%%%%%%%%%%%%%%%%%%%%%

%%%%%%%%%%%%%%%%%%%%%%%%%%%%%%%%%%%
% ANHÄNGE
%
% @stud: einzelne Anhänge bearbeiten und eigene Anhänge hier einfügen 
%        die nachfolgenden Zeilen deaktivieren, wenn keine Anhänge verwendet werden
%
\initializeAppendix

%%%%%%%%%%%%%%%%%%%%%%%%%%%%%%%%%%%

\singlespacing

%\ihead{}
%\printbibliography[title=\Literaturverzeichnis] 
\printbibliography 
\cleardoublepage

%\initializeBibliography
%%%%%%%%%%%%%%%%%%%%%%%%%%%%%%%%%%%

%%%%%%%%%%%%%%%%%%%%%%%%%%%%%%%%%%%
% INDEX
% @stud: ggf. Index auskommentieren, wenn nicht benötigt
%
\addcontentsline{toc}{chapter}{Index}
\printindex

\end{document}
